\documentclass[a4paper,12pt]{article}
\usepackage{color}
\usepackage{graphicx, gensymb}
\usepackage{amsfonts, cancel, amssymb}
\usepackage{amsmath, amsthm, array, esvect, esint}
\usepackage{typearea}
\usepackage[a4paper,left=2cm,right=2cm,top=2cm,bottom=3cm,bindingoffset=5mm]{geometry}
\newcommand\numberthis{\addtocounter{equation}{1}\tag{\theequation}}
\newcommand{\bra}{}
\newcommand{\kom}{}
\newcommand{\brak}{}
\newcommand{\braket}{}
\newcommand{\intx}{}
\newcommand{\intr}{}
\newcommand{\kla}{}
\newcommand{\klam}{}
\newcommand{\klammer}{}
\newcommand{\oper}{}
\newcommand{\tecxt}{}
\newcommand{\Bra}{}
\newcommand{\Ket}{}

\begin{document}

\title{09 Spektraltheorem}
\author{Joachim Schwardt}
\date{\today}
\maketitle

\section{Spektraltheorem für kompakte Operatoren}
a)
\begin{proof}
Sei $T\phi=0$. Dann ist natürlich auch $T^\ast T\phi=0$, also $\ker T\subseteq\ker T^\ast T$. Außerdem gilt für den umgekehrten Fall $0=\brak\langle {\phi}| {T^\ast T\phi}\rangle=\|T\phi\|^2$ und somit $\ker T^\ast T\subseteq\ker T$, woraus die Behauptung folgt.
\end{proof}
b)
\begin{proof}
Die Eigenvektoren selbstadjungierter Operatoren sind orthogonal und die Eigenwerte reell. Zur Erinnerung: 
\begin{align*}
\lambda\brak\langle {\lambda}| {\lambda}\rangle&=\brak\langle {\lambda}| {T\lambda}\rangle=\brak\langle {T\lambda}| {\lambda}\rangle=\lambda^\ast\brak\langle {\lambda}| {\lambda}\rangle\\
\lambda\brak\langle {\mu}| {\lambda }\rangle&=\brak\langle {\mu }| {T\lambda}\rangle=\brak\langle {T\mu }| {\lambda}\rangle=\mu^\ast\brak\langle {\mu }| {\lambda}\rangle\\
&\Rightarrow\ \ 0=\brak\langle {\mu }| {\lambda}\rangle(\mu-\lambda)
\end{align*}
Seien also $\Ket | {\mu } \rangle,\,\Ket | {\lambda} \rangle\in\ker\kla\left( {T^\ast T-\mu,\lambda I} \right)$. Dann folgt die Behauptung aus folgender Rechnung:
\begin{align*}
\brak\langle {T\mu }| {T\lambda}\rangle&=\brak\langle {\mu }| {T^\ast T\lambda}\rangle=\lambda\brak\langle {\mu }| {\lambda}\rangle\\
\brak\langle {T\mu }| {T\lambda}\rangle&=\brak\langle {T^\ast T\mu }| {\lambda}\rangle=\mu\brak\langle {\mu }| {\lambda}\rangle\\
\Rightarrow\ \ 0&=\brak\langle {\mu }| {\lambda}\rangle\kla\left( {\mu-\lambda} \right)
\end{align*}
\end{proof}
c)
\begin{proof}
Da $T^\ast T$ selbstadjungiert und kompakt ist, gibt es nach dem Spektraltheorem $\mu_j\subseteq\mathbb{R}$ und ein ONS $\phi_j \in\mathcal{H}$ so, dass $T^\ast T=\sum\mu_j\oper | {\phi_j} \rangle\langle {\phi_j} |$. Sei $\psi_j\subseteq\mathcal{H}$ ebenfalls ein ONS. Dann ist $T^\ast T=\sum\mu_j\oper | {\phi_j} \rangle\langle {\psi_j} \oper | {\psi_j} \rangle\langle {\phi_j} |=\sum\kla\left( {\sqrt{\mu_j}\oper | {\psi_j} \rangle\langle {\phi_j} |} \right)^\ast\kla\left( {\sqrt{\mu_j}\oper | {\psi_j} \rangle\langle {\phi_j} |} \right)$. Da wir ONS gewählt haben, kann man $T=\sum\sqrt{\mu_j}\oper | {\psi_j} \rangle\langle {\phi_j} |$ setzen, da bei der Betrachtung von $T^\ast T$ aus dem Produkt der Summe wieder eine Summe der Produkte wird.
\end{proof}
\section{Polynome selbstadjungierter Operatoren}
a)
\begin{proof}
Nachrechnen.
\end{proof}
b)
\section{Konvergenzen in der Normtopologie}
a)
\begin{proof}
Es gilt für die Konvergenzen, dass $\tecxt\text{Normtopologie}>\tecxt\text{starke Operatortopologie}>\tecxt\text{schwache Operatortopologie}$. Die hinreichenden Bedingungen sind hierbei leicht zu beweisen:\\
Sei $\|T_n-T\|\rightarrow 0$. Dann gilt für alle $\phi\in\mathcal{H}$, dass $\|T_n\phi-T\phi\|\le\|T_n-T\|\cdot\|\phi\|\rightarrow 0$.\\
Sei nun $\|T_n\phi-T\phi\|\rightarrow 0$. Dann gilt für alle $\phi,\psi\in\mathcal{H}$ mit der Cauchy-Schwartz Ungleichung, dass $|\brak\langle {\psi}| {(T_n-T)\phi}\rangle|\le\|\psi\|\cdot\|(T_n-T)\phi\|\rightarrow 0$.
\end{proof}
b)
\begin{proof}
Es ist $\|T_n\|=n$ für $\phi_j=e_j$. Also liegt keine Konvergenz in der Normtopologie vor. Allerdings ist $\|T_n\phi -T\phi\|=\|\sum_j^n\oper | {\phi_j} \rangle\langle {\phi_j} |\phi\rangle -\Ket | {\phi} \rangle\|\rightarrow 0$ gerade die Folgerung aus der Parseval'schen Gleichung.
\end{proof}

\end{document}